\documentclass[10pt]{article}
\usepackage[letterpaper,margin=0.86in]{geometry}
\usepackage{amsmath,amssymb,booktabs,xcolor,hyperref}
\pdfmapfile{+cm.map}\pdfmapfile{+cmextra.map}\pdfmapfile{+symbols.map}
\hypersetup{colorlinks=true,urlcolor=blue!50!black,pdftitle={A finite sieve for the four-ninths factorial window},pdfauthor={Jordan Pollard, with OpenAI Codex}}
\setlength{\parindent}{0pt}\setlength{\parskip}{5pt}\setlength{\emergencystretch}{2em}
\newcommand{\thetar}{\theta}
\begin{document}
{\LARGE\bfseries A finite sieve for the four-ninths factorial window\par}
\smallskip
Jordan Pollard\hfill November 27, 2026\\
{\small OpenAI Codex developed the argument and prepared the Lean proofs.}

\section*{1. The remaining counting estimate}
Put $c=4/9$, $h=\lfloor n^c\rfloor$ and $b=5/9+1/10000$. An intermediate-prime step for Erd\H{o}s problem 420 asks for a fixed $\eta>0$ such that, for every sufficiently large $n$,
\[
  \sum_{n<m\le n+h}\ \sum_{\substack{p\mid m\\n^c<p\le n^b}} v_p(m)\log p
  \ \le\ (5/9-\eta)h\log n .
\]
Here $p$ ranges over primes. The formal reduction gives $\eta=1/5000$ from 48 estimates that count each pair $(m,p)$ once. It also proves the bound needed for all repeated prime factors.

Divide each of the following three intervals into 16 equal subintervals. For each subinterval $(u,v]$, take $\thetar(v)$ from its row.
\begin{center}
\begin{tabular}{lll}
\toprule
Range of exponents & Exponent pair & $\thetar(v)$ \\
\midrule
$[4/9,77/162]$ & $(1/14,11/14)$ & $(4/7)/(17/42-5v/7)$ \\
$[77/162,29/54]$ & $(1/6,2/3)$ & $1/(19/54-v/2)$ \\
$[29/54,50009/90000]$ & $(1/2,1/2)$ & $12$ \\
\bottomrule
\end{tabular}
\end{center}
Let $C_n(u,v)$ count all pairs $n<m\le n+h$, $p\mid m$, $n^u<p\le n^v$. The remaining analytic input is that, eventually, in each of these 48 fixed bands,
\[
 C_n(u,v)\le h\big((v-u)\thetar(v)+10^{-6}\big).
\]
Equation (19) in the 1996 paper cited below implies these inequalities on paper. Cover a band by at most $(v-u)\log n+1$ multiplicative intervals of ratio $e$. Its bound for each interval is uniform, $\theta$ is nondecreasing on each of the three ranges, and $n^c/h\to1$. The tolerance absorbs the endpoint and asymptotic errors.

\textbf{First powers.} Each occurrence in a band contributes at most $v\log n$. The kernel checks the exact rational comparison
\[
 c+\sum_{(u,v]}v\big((v-u)\thetar(v)+10^{-6}\big)<2499/2500 .
\]
This gives a saving of $1/2500$ before accounting for multiplicities.

\textbf{Repeated powers.} Map each pair with $p>h$, $p^2\mid m$ to $q=m/p^2$. For fixed positive $q$, consecutive values of $qp^2$ in this range are more than $h$ apart, so the map is injective. There are at most $\lfloor(n+h)/(h+1)^2\rfloor\le 2n^{1/9}$ such pairs for large $n$. Also $(h+1)^3>n+h$, so each contributes just one extra copy of $\log p$. The extra cost is eventually at most $h\log n/5000$. This leaves the claimed $1/5000$ saving.

\newpage
\section*{2. The integral and the finite upper sieve}
The finite argument on page 1 uses rational arithmetic throughout. The original continuous calculation has a larger margin. For the three full ranges in the table let
\[
 I=\sum_{[a,b']}\int_a^{b'}u\,\thetar(u)\,du .
\]
A separate Lean proof verifies
\[
                       4/9+I<0.993755.
\]
For either of the first two rows write $\thetar(u)=K/(A-Bu)$. The positive denominator on the full interval gives the identity
\[
 \int_a^{b'}\frac{Ku}{A-Bu}\,du=
 \frac{KA}{B^2}\log\!\frac{A-Ba}{A-Bb'}-\frac{K}{B}(b'-a).
\]
The two logarithm arguments are $99/74$ and $37/27$. Setting $x=25/173$ and $5/32$, respectively, gives $(1+x)/(1-x)$ in each case. Lean bounds the logs by the first four odd terms in the expansion of $\tfrac12\log((1+x)/(1-x))$ with remainder at most $x^9/(1-x^2)$. The last row has integrand $12u$. The calculation certifies a margin greater than $0.00624$ in the limiting sieve coefficient.

\textbf{Finite sieve inequality.} For a finite block $S$ of positive moduli above $z$, fix a finite set $D\subseteq[1,z]$ with $1\in D$, real weights $w_1=1$, and nonnegative multiplicities $N(m)$. Then
\[
  \sum_{p\in S\,\mathrm{prime}}N(p)
  \ \le\ \sum_{d,e\in D}w_dw_e
         \sum_{\substack{m\in S\\{[d,e]}\mid m}}N(m),
\]
where $[d,e]$ is their least common multiple. Indeed, the squared divisor sum
$(\sum_{d\in D,\ d\mid m}w_d)^2$ equals one at each prime in the block and is nonnegative elsewhere. Multiply by $N(m)$, sum over $S$, and expand. The Lean proof permits signed weights and retains the signs in the expanded sum.

For the factorial problem one takes $N(m)=\lfloor(n+h)/m\rfloor-\lfloor n/m\rfloor$. The outstanding formal task is to control the sums of these multiplicities with the level and exponent pairs in the table. The finite inequality supplies the starting point. The numerical and prime-power proofs supply the final steps.

\medskip
{\small\textbf{Reference.} P. Erd\H{o}s, S. W. Graham, A. Ivi\'c and C. Pomerance, \emph{On the number of divisors of $n!$}, Analytic Number Theory, vol. 1, Progress in Mathematics 138 (1996), 337--355. Lemma 3 and equations (14)--(19). \href{https://math.dartmouth.edu/~carlp/factorial.pdf}{Author's PDF}.}
\end{document}
